\section{Introduction et contexte général}
\subsection{Les enjeux de la recherche reproductible}
\begin{frame}
\frametitle{Reproductibilité : une définition générale}
On dispose :
\begin{itemize}
  \item d'une publication,
  \item du code et du jeu de données associés 
\end{itemize}

~\\
On voudrait :  
\begin{itemize}
  \item changer un paramètre et reproduire une figure
  \item ajouter une contribution au code pour étendre ses fonctionnalités
\end{itemize}

~\\
Qui ? 
  \begin{itemize}
      \item le chercheur lui-même
      \item des membres de son équipe
      \item d'autres chercheurs
  \end{itemize}
  \begin{itemize}  
       \item[\ding{220}] \textcolor{maincolor}{maintenant ou plus tard}
  \end{itemize}

\end{frame}


%%---------------------------------------------------------------
\begin{frame}
\frametitle{Motivations}

%\setbeamercolor{block title}{bg=myblue3, fg=white}
%\setbeamercolor{block body}{bg=myblue3!15}
\setbeamercolor{block body}{bg=white}

\begin{block}{L'intérêt d'assurer la reproductibilité des résultats ?}
\begin{itemize}
\item Pour gagner du temps \textcolor{maincolor}{pour vous} : \\
\end{itemize}

\mycolor{Exemples :}

  \begin{itemize}
\item[\ding{220}] vous avez une suggestion d'un lecteur, d'un collègue pour améliorer un point
  \item[\ding{220}] vous reprenez votre code 5 ans après ...
  \item[\ding{220}] Un étudiant en thèse veut comparer ses résultats avec ceux d'un étudiant précédent \\  
  \end{itemize}

 \begin{itemize}
\item Pour faciliter la réutilisation de votre travail, même localement \\   
\end{itemize}
 
\end{block}

\end{frame}



%%---------------------------------------------------------------
\begin{frame}
\frametitle{Bénéfices pour la recherche}

%\setbeamercolor{block title}{bg=myblue3, fg=white}
%\setbeamercolor{block body}{bg=myblue3!15}
\setbeamercolor{block body}{bg=white}

\begin{itemize}
\item Donner confiance au lecteur/reviewer \\  
\item Augmenter la rapidité de diffusion des savoirs
\item Augmenter l'efficacité de la recherche en réutilisant des choses déjà faites
\item Augmenter la qualité de la recherche : un bug peut être signalé
\item Favorise les collaborations : des contributions peuvent être apportées, ...\\
\end{itemize}
~\\
\mycolor{\ding{220}  Augmenter la confiance dans les produits de la recherche, pour soi, pour les collègues et pour le grand public en général}

\end{frame}


%%---------------------------------------------------------------
\begin{frame}
  \frametitle{En pratique, pour le lecteur d'une publication}
  \vspace{-0.5cm}
\begin{columns}
  \begin{column}{0.5\textwidth}
De quoi a-t-on besoin ?
  \begin{itemize}
      \item article, disponible sur une archive
      \item code, disponible grâce à un lien web, un identifiant pérenne
      \item jeu de données, disponible sur le web (une page ou un entrepôt)
\end{itemize} 
  \end{column}
%%%
\begin{column}{0.5\textwidth}
\includegraphics[width=5cm]{./figures/dessin_RR.png}
\end{column}
%
\end{columns}
~\\
\mycolor{\ding{220}} Les principes de la \mycolor{Science Ouverte} offre un cadre idéal pour assurer la reproductibilité d'un travail.
\end{frame}


%%--------------------------------------------------------------
\begin{frame}
\frametitle{Contexte légal : science ouverte}

\begin{block}{Cadre légal}
\begin{itemize}
\item Science Ouverte : loi pour une république numérique, plan  \\
\item Loi pour l'intégrité scientifique \\   
\end{itemize}

\mycolor{Obligations}

\begin{itemize}
\item[\ding{220}] on doit ouvrir ses données, ses codes et ses publications 
\item[\ding{220}] 
\item[\ding{220}]  
\end{itemize}
\end{block}

\end{frame}

%%--------------------------------------------------------------
\begin{frame}
\frametitle{Contexte général : science ouverte}

\begin{block}{Organismes français}
\begin{itemize}
\item[\ding{220}] COSO :
  \begin{itemize}
  \item assure la mise ne \oe uvre de la politique de Science Ouverte
  \item édite des guides de bonnes pratiques : passeports pour la science ouverte
\end{itemize}    
\item[\ding{220}] Ateliers de la donnée : gestion des données de la recherche 
\item[\ding{220}]  
\end{itemize}
\end{block}


\begin{block}{Organismes à l'international}
\begin{itemize}
\item[\ding{220}] Politiques et projets européens autour de l'Open Science 
  
\item[\ding{220}] Ateliers de la donnée : gestion des données de la recherche 
\item[\ding{220}]  
\end{itemize}
\end{block}

\end{frame}

%%--------------------------------------------------------------
\begin{frame}
\frametitle{Contexte général : intégrité scientifique}

\begin{block}{}
\begin{itemize}
\item Ouvrir ses données, ses codes, ses artciles sauf impossibilité   \\
\item Les référencer correctement \\   
\end{itemize}

\mycolor{Science ouverte == socle de la reproductibilité}
\end{block}

\end{frame}



\subsection{Contexte historique}
%%---------------------------------------------------------------
\begin{frame}
\frametitle{Historique}
\vspace{-0.4cm}
%\begin{block}{Diverses études sur la reproductibilité des résultats : meta-science}
  \footnotesize
\begin{itemize}\footnotesize
\item 2016 : Questionnaire en ligne envoyé à 1576 chercheurs de disciplines très différentes
  \begin{itemize}\footnotesize
  \item 70 $\%$ ont échoué à reproduire un travail effectué par un autre chercheur
  \item + de 50 $\%$ ont échoué à reproduire leur propre expérience \\
 \publi{Baker M. 2016, 1,500 scientists lift the lid on reproducibility, Nature 533:452–454}
\end{itemize}
\end{itemize}
\vspace{-0.25cm}
\begin{columns}
  \begin{column}{0.5\textwidth}
\hspace{2cm}\includegraphics[width=4cm]{./figures/Nature_2.png}
  \end{column}
%%%
\begin{column}{0.5\textwidth}
\includegraphics[width=4cm]{./figures/Nature_1.png}
\end{column}
%
\end{columns}
%\end{block}
\end{frame}

%%---------------------------------------------------------------
\begin{frame}
  \frametitle{Causes du manque de reproductibilité}
\setbeamercolor{block body}{bg=white}

\begin{itemize}
      \item fraude, erreur humaine  \\
      \item manque de sensibilisation / formation / support sur ces enjeux
      \item course à la publication : manque de temps
      \item ...
\end{itemize}

\end{frame}

%%---------------------------------------------------------------
\begin{frame}
\frametitle{Historique}
\vspace{-0.4cm}
\begin{block}{Diverses études sur la reproductibilité des résultats : méta-recherche}
\footnotesize
\begin{itemize}
\footnotesize  
\item Economie
  \begin{itemize}
   \footnotesize 
  \item Reproductibilité de 60 articles (données + code + préanalyse disponibles)
  \item 1/3 papiers reproduits, +10$\%$ après contact avec les auteurs
    \publi{Chang, Andrew C., and Phillip Li. 2017. "A Preanalysis Plan to Replicate Sixty Economics Research Papers That Worked Half of the Time." American Economic Review, 107 (5): 60–64. DOI: 10.1257/aer.p20171034}
  \end{itemize}
  \footnotesize
  \vspace{-0.3cm}
\item Psychologie appliquée
  \begin{itemize}
    \footnotesize
  \item Reproductibilité de 35 articles publiés entre 2015 et 2018 dans la revue \textit{Cognition} (données disponibles)
  \item 11 articles reproduits, 11 après contacts avec les auteurs : 63$\%$ max \\
    \publi{Estimating the Prevalence of Transparency and Reproducibility-Related Research Practices in Psychology (2014–2017), Tom E. Hardwicke, Robert T. Thibault, […], and John P. A. Ioannidis, Volume 17, Issue 1, https://doi.org/10.1177/1745691620979806 }
  \end{itemize}
  \footnotesize
  \vspace{-0.3cm}
\item Ten Years Reproducibility Challenge Oct 11, 2019
  \begin{itemize}
   \footnotesize 
  \item Objectif : refaire tourner son code publié 10 ans après 
  \item Difficultés : retrouver le code + env. logiciel \\
   \publi{Challenge to scientists: does your ten-year-old code still run? Nature vol. 584 iss. 7822, 2020}
  \end{itemize}  
\end{itemize}
\end{block}
\end{frame}






%%---------------------------------------------------------------
\begin{frame}
  \frametitle{Contexte international}
  \vspace{-0.5cm}
  \begin{itemize}
  \item Beaucoup de communautés différentes se sont emparées du sujet
  \item Création de plusieurs réseaux en Europe et dans le monde  : UKRN en 2019,...
  \item Projets européens : horizon europe (programme européen
pour la recherche et l'innovation) --> projets autour de la metascience et de la reproductibilité (OSIRIS, TRUSTparency, iRISE, ...) 
  \end{itemize}
  \vspace{0.5cm} 
  \centering
\includegraphics[height=2.5cm]{./figures/world-map.png}\\
  
\end{frame}

%%---------------------------------------------------------------
\begin{frame}
  \frametitle{Contexte français}
  \vspace{-0.5cm}  

\hspace{10cm}\includegraphics[height=1.5cm]{./figures/FRRN_logo_crop.png}\\
%\hspace{10cm}\html{https://www.recherche-reproductible.fr}
  \begin{itemize}
  \item 2022 : création du réseau français
    \item fort soutien du ministère
  \end{itemize}
  
  \begin{itemize}
  \item adhésion : abonnement à la liste de diffusion
  \item Vous pouvez contribuer :
      \begin{itemize}
      \item proposer des wébinaires
      \item contribuer au GT ``Logiciels'' (rédiger des fiches de bonnes pratique)
      \item alimenter la base de données bibliograhique
      \item proposer une synthèse d'un article en lien avec la reproductibilité 
      \end{itemize}
  \end{itemize}

\hspace{10cm}\html{https://www.recherche-reproductible.fr}
\end{frame}


\begin{frame}
  \frametitle{Les revues évoluent également}

\begin{block}{Publications}
 \begin{itemize}\footnotesize
 \item Revues qui ont des pré-requis (variable en fonction des disciplines)
   \begin{itemize}\footnotesize
   \item IPOL (machine learning), Computo (statistique), ReScience C (numérique) 
   \end{itemize}  
 \item Badges de l'ACM (Association for Computing Machinery) : \html{https://www.acm.org/publications/policies/artifact-review-and-badging-current}
 \item CASCAD en économie (\html{https://www.cascad.tech/})
 \end{itemize} 
\end{block}
%
\end{frame}



\subsection{Terminologie}
%%---------------------------------------------------------------
\begin{frame}
\frametitle{Définitions / thématiques}
\vspace{-0.2cm}
\begin{block}{Différents type de reproductibilité en fonction des contextes de recherche}
\footnotesize
\begin{itemize}  
\item \mycolor{statistique} : trace de la raison qui fait qu'un résultat est significatif au sens statistique
    \begin{itemize}\footnotesize
    \item justification du choix du test statistique
    \item choix des paramètres du modèle
    \item des valeurs de seuil, de la taille des échantillons
    \end{itemize}  
\item \mycolor{empirique} : trace des conditions dans lesquelles l’expérience a eu lieu
        \begin{itemize}\footnotesize
        \item paramètres : le PH, la température…,
        \item ingrédients qui ont composé cette expérience (les espèces chimiques, ...)
        \item protocole suivi pour la réaliser (ordre des étapes, matériel utilisé…).
        \end{itemize}  
        \mycolor{\ding{220}} La reproductibilité empirique est souvent difficile à atteindre parce que certaines expériences requièrent une précision importante dans la reproduction des conditions expérimentales (dosage exact des composants…).
 \end{itemize}         
\end{block}
\end{frame}




\begin{frame}
\frametitle{Définitions / thématiques}
\vspace{-0.2cm}
\begin{block}{Différents type de reproductibilité en fonction des contextes de recherche}
\footnotesize
\begin{itemize}

      \item \mycolor{computationnelle} : trace de la production d'un résultat sur ordinateur
    \begin{itemize}\footnotesize
    \item codes informatiques utilisés + leur version + enchainement
    \item environnement logiciel, type de machine, logiciels installés utilisés
    \item les jeux de données
    \end{itemize}
    \mycolor{\ding{220}} Pour la reproductibilité computationnelle : il existe plein d'outils informatiques qui facilitent grandement le travail !
  \end{itemize}

\end{block}
\end{frame}


%%---------------------------------------------------------------
\begin{frame}
  \frametitle{Définitions qui font consensus pour le calcul numérique}
  \vspace{-0.5cm}
  \begin{itemize}
      \item \mycolor{Reproductibilité} : en utilisant le même code (protocole, ..) et les mêmes données,
      on aboutit à des \textit{conclusions/résultats/..} \mycolor{identiques}
      \item \mycolor{Réplicabilité} : en prenant un jeu de données différents et/ou
      une implémentation différente de la méthode, on aboutit à des \textit{conclusions/résultats/..} \mycolor{équivalents}
  \end{itemize}
  \vspace{-0.2cm}
\centering
\begin{figure}
\includegraphics[width=6cm]{./figures/reproducible-matrix.jpg}
\end{figure}
\publi{The turing way: https://book.the-turing-way.org/}

\vspace{0.3cm}
Chaque domaine de recherche a ses contraintes propres !
\end{frame}




%%## Partie II : reproductibilité computationnelle et HPC
%%---------------------------------------------------------------
\section{Reproductibilité computationnelle}
\subsection{Problématique}
\begin{frame}
  \frametitle{Exemples de problèmes de reproductibilité en calcul scientifique}
  %\framesubtitle{version + options du compilateur}

\begin{tabular}{lr}

\end{tabular}

\vspace{-0.5cm}  
   
\centering
\hspace{2cm}
\begin{tabular}{lr}
Le code & Les dépendances \\  
\includegraphics[height=3cm]{./figures/make-arbre-dependances.png}
& \includegraphics[height=3.5cm]{./figures/maven-tree-plugin.png}\\
\end{tabular}
\end{frame}

%%---------------------------------------------------------------
\begin{frame}
  \frametitle{Exemples de problèmes de reproductibilité en calcul scientifique}
  \framesubtitle{version + options de python + dépendances}
  Version de python :\\
  En Python 3.6, les valeurs par défaut de rel$\_$tol et abs$\_$tol ont été ajustées : 
\begin{tabular}{|c|c|}
\hline
python3.5 & python3.6 \\
import math & import math \\
math.isclose(1.0, 1.0000001) :: False & math.isclose(1.0, 1.0000001) :: True \\
\hline
\end{tabular}

\end{frame}

%%---------------------------------------------------------------
\begin{frame}[fragile]
  \frametitle{Exemples de problèmes de reproductibilité en calcul scientifique}
  %\framesubtitle{version + options de python + dépendances}
  \mycolor{Dépendances} : changement de syntaxe \\
 \vspace{-0.2cm}  
 \begin{minted}[bgcolor=gray!10,fontsize=\scriptsize]{python}
import numpy
arr = np.array([[1, 2], [1, 2]])   
\end{minted} 
 \begin{tabular}{|c|c|}
\hline
NumPy 1.13 & NumPy >= 1.13 \\
np.unique(arr) & np.unique(arr, axis=0) \\
Résultat : array([1, 2])  & Résultat : array([[1, 2]]) \\
\hline
 \end{tabular}


 \vspace{0.25cm} 


\begin{tabular}{|c|c|}
\hline
Pandas < 0.20 & Pandas >= 0.20 \\
df.sort('column$\_$name') &df.sort$\_$values('column$\_$name') \\
\hline
\end{tabular}
\end{frame}



%%---------------------------------------------------------------
\begin{frame}[fragile]
  \frametitle{Exemples de problèmes de reproductibilité en calcul scientifique}
  %\framesubtitle{version + options de python + dépendances}
  \mycolor{Dépendances} : résultat d'une fonction 
\vspace{0.15cm}
\begin{minted}[bgcolor=gray!10,fontsize=\scriptsize]{python}
import random
random.seed(42)
data = [1, 2, 3, 4, 5]
random.shuffle(data)
print(data)
\end{minted}

\begin{tabular}{|c|c|}
\hline
random 3.2 & random 3.3+ \\
Résultat : [4, 2, 3, 5, 1] & Résultat : [1, 5, 2, 4, 3] \\
\hline
\end{tabular}
\end{frame}




%%---------------------------------------------------------------
\begin{frame}[fragile]

\frametitle{Exemples de problèmes de reproductibilité en calcul scientifique}
\framesubtitle{Parallélisation avec OpenMP}
\vspace{-0.25cm}
  $$H_n = \sum_{k=1}^n 1/k = ln(n) + \gamma +o(1)$$
\begin{itemize}
\item Intel(R) Xeon(R) Gold 6130 CPU @ 2.10GHz 
\item gfortran 14 -fopenmp 
\end{itemize}

\begin{minted}[bgcolor=gray!10,fontsize=\scriptsize]{fortran}
Gamma = 0.5772156649015328606
!$OMP parallel do private(i) reduction(+ : sum)
do i=1,n
   sum = sum + 1.0/i
enddo
!$OMP end parallel do 
p = abs(sum - log(float(n)) - Gamma)
\end{minted}
\end{frame}



%%---------------------------------------------------------------
\begin{frame}[fragile]

\frametitle{Problèmes de reproductibilité : parallélisation avec OpenMP}
%\framesubtitle{Parallélisation avec OpenMP}
\vspace{-0.25cm}
\footnotesize
n= 1000000 \\

\begin{tabular}{|c|c|c|}
\hline  
Calcul & séquentiel & OpenMP 1 thread  \\
\hline
sum &14.392726788474306  &14.392726788474306  \\
p   &3.6711810480483109E-007 &3.6711810480483109E-007 \\
\hline
\end{tabular}

\begin{tabular}{|c|c|}
\hline  
 OpenMP 4 threads & OpenMP 4 threads -Ofast \\
\hline
14.392726\maincolor{788474306} &14.392726\maincolor{222394913} \\
 3.6711810480483109E-007 &1.9896128833352122E-007 \\
\hline
\end{tabular}



\vspace{0.2cm}
\footnotesize
n= 100000000 \\

\begin{tabular}{|c|c|c|}
\hline  
Calcul & séquentiel & OpenMP 1 thread  \\
\hline
sum & 18.9978964\maincolor{82991536}  & 18.9978964\maincolor{82991536}  \\
p   & 1.8830363046618004E-007  & 1.8830363046618004E-007 \\
\hline
\end{tabular}

\begin{tabular}{|c|c|c|}
\hline  
 OpenMP 4 threads & OpenMP 4 threads -Ofast \\
\hline
 18.9978964\maincolor{79461851} & 18.99789\maincolor{5850727968}  \\
 1.9183331545491455E-007 & 8.2056719818979218E-007\\
\hline
\end{tabular}
\end{frame}






%%---------------------------------------------------------------
\begin{frame}[fragile]

\frametitle{Problèmes de reproductibilité en calcul scientifique : calcul sur GPU}
%\framesubtitle{Calcul sur GPU}
\vspace{-0.25cm}
  $$H_n = \sum_{k=1}^n 1/k = ln(n) + \gamma +o(1)$$
\begin{itemize}
\item Intel(R) Xeon(R) Gold 6130 CPU @ 2.10GHz 
\item 1 GPU nvidia V100
\item nvfortran 22.3-0 64-bit target on x86-64 Linux
\end{itemize} 
\begin{minted}[bgcolor=gray!10,fontsize=\scriptsize]{fortran}
Gamma = 0.5772156649015328606
!$acc kernels copy(sum)
devicetype = acc_get_device_type()
ngpus = acc_get_num_devices(devicetype)
do i=1,n
   sum = sum + 1.0/i
enddo
!$acc end kernels
p = abs(sum - log(float(n)) - Gamma)
\end{minted}
\end{frame}
%%---------------------------------------------------------------
\begin{frame}[fragile]

  \frametitle{Exemples de problèmes de reproductibilité en calcul scientifique}
  %\framesubtitle{calcul parallèle / GPU}

\footnotesize
n= 1000000 \\

\begin{tabular}{|c|c|c|c|}
\hline  
Calcul & CPU & CPU option fast & GPU \\
\hline
sum &14.39272678847431  & 14.39272678847431& 14.39272678847431 \\
p   & 3.6711810480483109E-007 & 3.6711810480483109E-007 & 3.6711810480483109E-007 \\
\hline
\end{tabular}

\vspace{0.5cm}
\footnotesize
n= 100000000 \\

\begin{tabular}{|c|c|c|c|}
\hline  
Calcul & CPU & CPU option fast & GPU \\
\hline
sum & 18.99789648\maincolor{299154} & 18.99789648\maincolor{681272} & 18.997896\maincolor{47946185} \\
p   & 1.8830363046618004E-007 & 1.8448244887281362E-007 & 1.9183331545491455E-007  \\
\hline
\end{tabular}
\end{frame}


\subsection{Les bonnes pratiques}

%%---------------------------------------------------------------
\begin{frame}
  \frametitle{Calcul scientifique : les bonnes pratiques}
  
\begin{itemize}
\item Gestion des données : utilisation et production 
  \item Gestion des codes : écriture, exécution et diffusion 
  \end{itemize}
\vspace{0.4cm}
\center
\mycolor{\ding{220} Pour chaque étape il y a des outils pour vous aider !}
\end{frame}

%%---------------------------------------------------------------
\begin{frame}
  \frametitle{Calcul scientifique : les bonnes pratiques}
%  \framesubtitle{Gestion des données}

  \begin{block}{Gestion des données d'entrée }
Dans une publication, vous indiquerez : 
  \begin{itemize}
  \item DOI + lien téléchargement
  \item Trace d'une conversion de format 
  \item Trace du traitement effectué
  \item La façon dont elles sont utilisées dans votre simulation 
  \end{itemize}
\end{block}

\end{frame}

%%---------------------------------------------------------------
\begin{frame}
  \frametitle{Calcul scientifique : les bonnes pratiques}
%  \framesubtitle{Gestion des données}

\begin{block}{Gestion des données de sortie : respect des principes FAIR }
  \begin{itemize}
  \item format standard
  \item licence
  \item DOI + entrepôt
  \item Description 
  \item Métadonnées 
\end{itemize}
\end{block}  
\end{frame}


%%---------------------------------------------------------------
\begin{frame}
  \frametitle{Calcul scientifique : les bonnes pratiques}
  \framesubtitle{Ecriture et exécution du code}

\begin{itemize}
\item Mettre en \oe uvre les bonnes pratiques de développement
\item Si possible, utiliser un \mycolor{gestionnaire d'environnement logiciel reproductible (Guix, Nix)}
\item Réfléchir à une licence : va réglementer la vie du logiciel 
\item Utiliser un \mycolor{Outil de versionnement} (forge logicielle)
\item Rédiger une \mycolor{documentation du code} \\
\end{itemize}

\vspace{0.3cm}

Penser à sauvegarder dans des fichiers adaptés : 
\begin{itemize}
\item Description des moyens de calcul (portable, cluster, ...)[README]
\item \mycolor{Environnement logiciel}
\item Jeux de paramètres 
\item \mycolor{Workflow} d'éxécution
\end{itemize}

\end{frame}



%%---------------------------------------------------------------
\begin{frame}
  \frametitle{Calcul scientifique : les bonnes pratiques}
  \framesubtitle{Diffusion du code}

  \begin{itemize}
  \item Mettre une license (licence, contributeurs, etc ..)
  \item Identifiant pérenne : HAL + SWH
    \includegraphics[height=3cm]{./figures/Git_Swh_Hal_3.png}
  \item Description de l'environnement logiciel 
  \item Fichier de workflow 
\end{itemize}

\end{frame}



\section{Conclusions}
%% III Les leviers d'actions 
%%---------------------------------------------------------------

\begin{frame}
  \frametitle{Personne n'est parfait !}

\begin{block}{Idée : changer ses pratiques en douceur \emoji{winking-face} }  

\begin{itemize}
  \item identifier les facteurs qui freinent l'accès à la reproductibilité
    \begin{itemize}
  \item[$\circ$] facteurs humains
  \item[$\circ$] facteurs techniques
  \item[$\circ$] facteurs structurels : incitation à la publication
 \end{itemize}   
 \item Identifier les points à améliorer 
 \item Comment ? il existe des outils 
\end{itemize}
\end{block}
%
\end{frame}

%%---------------------------------------------------------------
\begin{frame}
  \frametitle{Conclusions}
\begin{block}{En pratique }
\begin{itemize}
\item On perd un peu de temps à la mise en place du protocole 
\item On en gagne vraiment beaucoup \textit{a posteriori}
\item Il existe pas mal d'outils qui aident vraiment 
\end{itemize}
\end{block}

\begin{block}{Qui est gagnant ? }
\begin{itemize}
\item Vous : rigueur + gain de temps
\item Les reviewers : confiance
\item Vos collègues : facilite la transmission des codes, les collaborations  
\item Le grand public : confiance
\end{itemize}
\end{block}
\end{frame}


%%---------------------------------------------------------------
\begin{frame}
  \frametitle{Exercice}

\begin{block}{Mise en pratique }
\begin{itemize}
\item par groupe de 4 ou 5 
\item présentations
\item identifier une bonne pratique que vous avez 
\item réfléchissez à un élément que vous pourriez facilement améliorer
\item discussion 
\item un membre du groupe fait une restitution 
\end{itemize}
\end{block}
\end{frame}

%%% Local Variables:
%%% mode: latex
%%% TeX-master: "main"
%%% End:
